% ============================================================
% CAPITOLO K - LA MODULAZIONE AI TONI VICINI (VERSIONE ESTESA)
% ============================================================
\chapter{La modulazione ai toni vicini}
\label{chap:modulazione}

Nei capitoli precedenti abbiamo costruito il nostro edificio armonico all'interno di un'unica tonalità, imparando a muoverci tra le funzioni di tonica, sottodominante e dominante con fluidità e coerenza. Abbiamo classificato le sette specie di accordi di settima (Capitolo \ref{chap:sette_specie}), ognuna con la propria struttura e funzione privilegiata. Ora è il momento di imparare a viaggiare: la \textbf{modulazione} è l'arte di spostarsi da una tonalità a un'altra, aprendo nuove prospettive espressive e creando architetture musicali più ampie e complesse.

La modulazione non è un salto improvviso nel vuoto, ma un processo organico e logico, che sfrutta le relazioni naturali tra le tonalità e la capacità di certi accordi di avere un doppio significato armonico. In questo capitolo ci concentreremo sulle modulazioni più naturali e immediate: quelle verso i \textbf{toni vicini}, ovvero le tonalità che condividono la maggior parte delle note con la tonalità di partenza.

\section{Che cos'è una modulazione}
\label{sec:definizione_modulazione}

Prima di addentrarci nelle tecniche, è necessario definire con precisione cosa intendiamo per ``modulazione'' e distinguerla da altri fenomeni armonici.

\begin{definizione}[Modulazione]
	\label{def:modulazione}
	Una \textbf{modulazione} è il processo mediante il quale una composizione musicale abbandona stabilmente una tonalità di partenza (tonalità di impianto) per stabilirne una nuova. La modulazione comporta:
	\begin{enumerate}
		\item L'affermazione chiara della tonalità di partenza
		\item Un percorso armonico di transizione
		\item La conferma della nuova tonalità attraverso una cadenza
	\end{enumerate}
	
	Si distingue dal semplice \textbf{accenno} o \textbf{tonicizzazione}, in cui si crea momentaneamente l'illusione di una nuova tonalità senza abbandonare stabilmente quella di partenza.
\end{definizione}

La differenza tra modulazione vera e propria e tonicizzazione è sottile ma importante: nella tonicizzazione, la nuova tonalità viene ``visitata'' brevemente ma non confermata da una cadenza autentica; nella modulazione, invece, il cambiamento è strutturale e la nuova tonalità diventa il nuovo centro di gravità armonico.

\section{Le tre fasi del processo modulante}
\label{sec:fasi_modulazione}

Ogni modulazione ben costruita si articola in tre momenti distinti, ciascuno con una funzione specifica.

\begin{regola}[Le tre fasi della modulazione]
	\label{reg:tre_fasi_modulazione}
	Il processo modulante si sviluppa attraverso:
	
	\begin{enumerate}
		\item \textbf{Affermazione (o consolidamento):} La tonalità di partenza viene stabilita chiaramente attraverso progressioni caratteristiche e, preferibilmente, una cadenza. L'ascoltatore deve percepire senza ambiguità qual è il ``centro tonale'' da cui si parte.
		
		\item \textbf{Transizione (o modulazione propriamente detta):} Si introduce un \textbf{accordo pivot} (o accordo cardine), ovvero un accordo che appartiene simultaneamente a entrambe le tonalità. Questo accordo funge da ``ponte'' armonico, permettendo il passaggio dalla tonalità vecchia a quella nuova senza discontinuità percettive.
		
		\item \textbf{Conferma (o cadenza):} La nuova tonalità viene confermata attraverso una progressione cadenzale chiara, tipicamente una cadenza perfetta (V--I) o composta (II--V--I o IV--V--I). Questa fase è essenziale: senza una cadenza convincente nella nuova tonalità, abbiamo solo una tonicizzazione passeggera, non una modulazione stabile.
	\end{enumerate}
\end{regola}

Questo schema tripartito non è una regola rigida ma un modello ideale. Nelle composizioni reali, specialmente in quelle più brevi, la fase di affermazione può essere implicita (se la tonalità di partenza è già ben stabilita) e la transizione può essere molto rapida. Ciò che non può mai mancare è la \textbf{conferma} della nuova tonalità.

\section{I toni vicini: una geografia armonica}
\label{sec:toni_vicini}

Non tutte le tonalità sono ugualmente ``lontane'' l'una dall'altra. Alcune condividono molte note, altre poche o nessuna. Questa ``distanza'' armonica è visualizzabile attraverso il circolo delle quinte (cfr. Figura \ref{fig:circolo_quinte_migliorato_compatto}).

\begin{definizione}[Tonalità vicine]
	\label{def:tonalita_vicine}
	Due tonalità si dicono \textbf{vicine} (o adiacenti) se le loro armature di chiave differiscono per non più di \textbf{un'alterazione} (un diesis o un bemolle). Questa vicinanza si traduce nella condivisione di sei note su sette della scala diatonica.
	
	Dal punto di vista del circolo delle quinte, le tonalità vicine sono quelle immediatamente adiacenti (distanza di una quinta) o legate dal rapporto di relativa maggiore/minore.
\end{definizione}

\subsection{I cinque toni vicini di una tonalità maggiore}
\label{subsec:toni_vicini_maggiore}

\begin{regola}[Tonalità vicine al modo maggiore]
	\label{reg:toni_vicini_maggiore}
	Data una tonalità maggiore, i suoi cinque toni vicini sono:
	
	\begin{enumerate}
		\item La \textbf{relativa minore} (stessa armatura di chiave)
		\item La tonalità della \textbf{dominante} (una quinta sopra, un diesis in più o un bemolle in meno)
		\item La \textbf{relativa minore della dominante} (stessa armatura della dominante)
		\item La tonalità della \textbf{sottodominante} (una quinta sotto, un bemolle in più o un diesis in meno)
		\item La \textbf{relativa minore della sottodominante} (stessa armatura della sottodominante)
	\end{enumerate}
\end{regola}

\begin{esempio}[I cinque toni vicini di Do Maggiore]
	\label{ex:toni_vicini_do}
	Partendo da Do Maggiore (nessuna alterazione), i toni vicini sono:
	
	\begin{itemize}
		\item \textbf{La minore} (relativa, 0 alterazioni)
		\item \textbf{Sol Maggiore} (dominante, 1\#)
		\item \textbf{Mi minore} (relativa di Sol, 1\#)
		\item \textbf{Fa Maggiore} (sottodominante, 1$\flat$)
		\item \textbf{Re minore} (relativa di Fa, 1$\flat$)
	\end{itemize}
	
	Tutte queste tonalità condividono sei note su sette con Do Maggiore.
\end{esempio}

\subsection{I cinque toni vicini di una tonalità minore}
\label{subsec:toni_vicini_minore}

\begin{regola}[Tonalità vicine al modo minore]
	\label{reg:toni_vicini_minore}
	Data una tonalità minore, i suoi cinque toni vicini sono:
	
	\begin{enumerate}
		\item La \textbf{relativa maggiore} (stessa armatura di chiave della minore naturale)
		\item La tonalità del \textbf{III grado} (relativa maggiore, ma costruita sul III della scala minore naturale)
		\item La tonalità del \textbf{IV grado} (sottodominante)
		\item La tonalità del \textbf{V grado} (dominante)
		\item La tonalità del \textbf{VI grado} (sopradominante)
	\end{enumerate}
	
	Nota: Nel modo minore, la situazione è leggermente più complessa a causa delle varianti armoniche e melodiche della scala, ma il principio della vicinanza per alterazioni rimane valido.
\end{regola}

\begin{esempio}[I cinque toni vicini di La minore]
	\label{ex:toni_vicini_la_m}
	Partendo da La minore (nessuna alterazione in chiave), i toni vicini sono:
	
	\begin{itemize}
		\item \textbf{Do Maggiore} (relativa, 0 alterazioni)
		\item \textbf{Mi minore} (dominante, 1\#)
		\item \textbf{Re minore} (sottodominante, 1$\flat$)
		\item \textbf{Sol Maggiore} (VII della minore = sottodominante della relativa, 1\#)
		\item \textbf{Fa Maggiore} (VI della minore = sottodominante della relativa, 1$\flat$)
	\end{itemize}
\end{esempio}

\section{L'accordo pivot: il ponte tra due mondi}
\label{sec:accordo_pivot}

Il cuore della tecnica modulante ai toni vicini è l'uso intelligente degli \textbf{accordi pivot} (detti anche accordi cardine o di perno).

\begin{definizione}[Accordo pivot]
	\label{def:accordo_pivot}
	Un \textbf{accordo pivot} è un accordo che appartiene simultaneamente a due tonalità diverse, avendo una funzione armonica specifica in ciascuna di esse. Nella modulazione, l'accordo pivot viene dapprima percepito come parte della tonalità di partenza, poi ``reinterpretato'' come parte della tonalità di arrivo.
	
	La doppia identità dell'accordo pivot permette una transizione fluida e naturale, senza discontinuità percettive.
\end{definizione}

Non tutti gli accordi sono ugualmente efficaci come pivot. La scelta dipende da due fattori:

\begin{enumerate}
	\item \textbf{Condivisione delle note:} L'accordo deve essere costruibile con le note di entrambe le scale
	\item \textbf{Forza funzionale nella nuova tonalità:} Un accordo pivot che nella nuova tonalità ha funzione di II, IV o V grado è più efficace di uno che ha funzione di III o VI
\end{enumerate}

\subsection{Accordi pivot tra tonalità relative}
\label{subsec:pivot_relative}

La modulazione tra tonalità relative (Do Maggiore $\leftrightarrow$ La minore) è la più semplice, poiché le due tonalità condividono tutte le note della scala diatonica.

\begin{regola}[Accordi pivot tra relative]
	\label{reg:pivot_relative}
	Tra Do Maggiore e La minore, gli accordi pivot più efficaci sono:
	
	\begin{itemize}
		\item \textbf{Re minore:} II grado in Do M, IV grado in La m (funzione di sottodominante in entrambe)
		\item \textbf{Mi minore:} III grado in Do M, V grado in La m (da tonica a dominante)
		\item \textbf{Fa Maggiore:} IV grado in Do M, VI grado in La m
	\end{itemize}
	
	Il più efficace è tipicamente il II di Do M / IV di La m, perché ha funzione di sottodominante forte in entrambe le tonalità.
\end{regola}

\begin{esempio}[Modulazione Do Maggiore $\rightarrow$ La minore tramite IV di La m]
	\label{ex:mod_do_lam}
	Progressione che modula usando Re minore come pivot.

	\musica{
		\partilate
		{ \key do \major \time 4/4 }
		{ mi'2 fa'2 | sold2 la2 | la1 }
		{ do'2 do'2 | mi2 mi2 | mi1 }
		{ sol2 la2 | si,2 do2 | do1 }
		{ do2 fa,2 | mi,2 la,2 | la,1 }
		\lyricmode { 
			\markup{"Do M:"} "I"2 "IV (pivot)"2 | 
			\markup{"La m:"} "V"2 "I"2 | 
			"I"1 
		}
		\figuremode { s2 s2 | <\markup{\sharp}>2 s2 | s1 }
	}
	Il Fa maggiore (IV di Do M) viene reinterpretato come IV di La minore, preparando la cadenza V--I in La minore.
\end{esempio}

\subsection{Accordi pivot verso la dominante}
\label{subsec:pivot_dominante}

La modulazione dalla tonica alla sua dominante (Do M $\rightarrow$ Sol M) è altrettanto naturale.

\begin{regola}[Accordi pivot verso la dominante maggiore]
	\label{reg:pivot_dominante}
	Tra Do Maggiore e Sol Maggiore, gli accordi pivot più comuni sono:
	
	\begin{itemize}
		\item \textbf{Do Maggiore stesso:} I grado in Do M, IV grado in Sol M
		\item \textbf{Mi minore:} III grado in Do M, VI grado in Sol M
		\item \textbf{Sol Maggiore (con settima):} V$^7$ in Do M, I$^7$ in Sol M (meno comune come pivot stabile)
	\end{itemize}
	
	Il più efficace è Do Maggiore reinterpretato come IV di Sol, perché permette di precedere immediatamente la cadenza V--I nella nuova tonalità.
\end{regola}

\begin{esempio}[Modulazione Do Maggiore $\rightarrow$ Sol Maggiore tramite pivot I = IV]
	\label{ex:mod_do_sol}

	\musica{
		\partilate
		{ \key do \major \time 4/4 }
		{ mi'2 mi'2 | fad'2 sol'2 | sol'1 }
		{ do'2 do'2 | do'2 si2 | si1 }
		{ sol2 sol2 | la2 sol2 | sol1 }
		{ do2 do2 | re2 sol,2 | sol,1 }
		\lyricmode { 
			\markup{"Do M:"} "I (pivot)"2 "I"2 | 
			\markup{"Sol M:"} "IV"2 "V"2 | 
			"I"1 
		}
		\figuremode { s2 s2 | s2 s2 | s1 }
	}
	L'accordo di Do Maggiore, percepito inizialmente come tonica, viene reinterpretato come sottodominante di Sol Maggiore.
\end{esempio}

\subsection{Accordi pivot verso la sottodominante}
\label{subsec:pivot_sottodominante}

\begin{regola}[Accordi pivot verso la sottodominante maggiore]
	\label{reg:pivot_sottodominante}
	Tra Do Maggiore e Fa Maggiore:
	
	\begin{itemize}
		\item \textbf{Fa Maggiore stesso:} IV in Do M, I in Fa M
		\item \textbf{Re minore:} II in Do M, VI in Fa M
		\item \textbf{Do Maggiore:} I in Do M, V in Fa M
	\end{itemize}
	
	Il più efficace è Do Maggiore reinterpretato come dominante di Fa, che porta direttamente alla cadenza.
\end{regola}

\begin{esempio}[Modulazione Do Maggiore $\rightarrow$ Fa Maggiore tramite pivot I = V]
	\label{ex:mod_do_fa}
	
	\musica{
		\partilate
		{ \key do \major \time 4/4 }
		{ mi'2 mi'2 | fa'2 mi'2 | fa'1 }
		{ do'2 sib2 | la2 do'2 | do'1 }
		{ sol2 sol2 | do2 sib2 | la1 }
		{ do2 do2 | fa,2 fa,2 | fa,1 }
		\lyricmode { 
			\markup{"Do M:"} "I"2 "I (pivot)"2 | 
			\markup{"Fa M:"} "V"2 "V"2 | 
			"I"1 
		}
		\figuremode { s2 s2 | s2 s2 | s1 }
	}
\end{esempio}

\section{Le sette specie come strumenti modulanti}
\label{sec:specie_modulazione}

La conoscenza delle sette specie di accordi di settima (Capitolo \ref{chap:sette_specie}) apre possibilità modulanti molto più ricche e sofisticate rispetto all'uso delle sole triadi.

\subsection{La seconda specie: il pivot per eccellenza}
\label{subsec:seconda_specie_pivot}

Come abbiamo visto nel Capitolo \ref{chap:vii7_ii7}, la settima semidiminuita (seconda specie) ha una natura doppia: può essere VII$\oslash^7$ nel modo maggiore o II$\oslash^7$ nel modo minore.

\begin{regola}[La seconda specie come accordo pivot]
	\label{reg:seconda_specie_pivot}
	L'accordo di settima semidiminuita è il pivot ideale tra tonalità relative:
	
	\begin{itemize}
		\item \textbf{VII$\oslash^7$ di Do Maggiore = II$\oslash^7$ di La minore}
		\item Stesso accordo (Si--Re--Fa--La), due funzioni: dominante in Do M, sottodominante in La m
	\end{itemize}
	
	Questa ambiguità funzionale permette modulazioni estremamente fluide tra tonalità relative.
\end{regola}

\begin{esempio}[Modulazione con VII$\oslash^7$ = II$\oslash^7$ come pivot]
	\label{ex:mod_seconda_specie}

	\musica{
		\partilate
		{ \key do \major \time 4/4 }
		{ mi'2 mi'2 | la'2 sold'2 | la'1 }
		{ do'2 re'2 | fa'2 mi'2 | mi'1 }
		{ sol2 si2 | do'2 si2 | do'1 }
		{ do2 si,2 | re2 mi2 | la,1 }
		\lyricmode { 
			\markup{"Do M:"} "I"2 "VII$\oslash^7$ (pivot)"2 | 
			\markup{"La m:"} "II$\oslash^7$"2 "V"2 | 
			"I"1 
		}
		\figuremode { s2 <6 5>2 | <6 5>2 s2 | s1 }
	}
	L'accordo Si--Re--Fa--La funge da cerniera: viene percepito come VII$\oslash^7$ (funzione di dominante) in Do Maggiore, poi reinterpretato come II$\oslash^7$ (funzione di sottodominante) in La minore.
\end{esempio}

\subsection{La terza specie: versatilità modulante}
\label{subsec:terza_specie_pivot}

La settima minore (terza specie) appare su molti gradi di diverse scale, rendendola estremamente versatile.

\begin{regola}[La terza specie come pivot tra tonalità maggiori]
	\label{reg:terza_specie_pivot}
	Esempi di accordi di terza specie usabili come pivot:
	
	\begin{itemize}
		\item \textbf{Re--Fa--La--Do:} II$^7$ in Do M = VI$^7$ in Fa M = VII$^7$ in Mi m
		\item \textbf{Mi--Sol--Si--Re:} III$^7$ in Do M = VI$^7$ in Sol M = II$^7$ in Re M
	\end{itemize}
	
	La terza specie permette di modulare non solo ai toni vicini immediati, ma anche di concatenare più modulazioni successive.
\end{regola}

\begin{esempio}[Modulazione Do M $\rightarrow$ Fa M tramite II$^7$ = VI$^7$]
	\label{ex:mod_terza_specie}
	
	\musica{
		\partilate
		{ \key do \major \time 4/4 }
		{ mi'2 do'2 | sib2 la2 | la1 }
		{ do'2 la2 | fa2 fa2 | fa1 }
		{ sol2 re2 | re2 do2 | do1 }
		{ do2 re,2 | sol,2 fa,2 | fa,1 }
		\lyricmode { 
			\markup{"Do M:"} "I"2 "II$^7$ (pivot)"2 | 
			\markup{"Fa M:"} "VI$^7$"2 "V"2 | 
			"I"1 
		}
		\figuremode { s2 <7>2 | <7>2 s2 | s1 }
	}
\end{esempio}

\subsection{La quinta specie: modulazioni audaci}
\label{subsec:quinta_specie_modulazione}

Come discusso nella Sezione \ref{sec:quinta_specie_dettaglio}, la settima diminuita ha una struttura simmetrica che le permette di risolvere su quattro toniche diverse.

\begin{regola}[La quinta specie per modulazioni multiple]
	\label{reg:quinta_specie_modulazioni}
	Un singolo accordo di settima diminuita può essere reinterpretato come sensibile di quattro tonalità diverse (distanti una terza minore l'una dall'altra). Questo apre la possibilità di modulazioni:
	
	\begin{itemize}
		\item Ai toni vicini (quando la nuova tonica dista una terza minore)
		\item Ai toni lontani (argomento che svilupperemo in un capitolo successivo)
	\end{itemize}
	
	Attenzione: L'uso della quinta specie per modulazioni ai toni lontani esce dall'ambito di questo capitolo e verrà trattato separatamente.
\end{regola}

Per ora, ci limitiamo a osservare che la quinta specie è uno strumento potentissimo, ma richiede grande attenzione nella preparazione e nella conferma della nuova tonalità per non creare ambiguità percettive.

\section{La tecnica della modulazione: consigli pratici}
\label{sec:tecnica_modulazione}

\subsection{Scegliere il pivot giusto}
\label{subsec:scelta_pivot}

\begin{regola}[Criteri per la scelta del pivot]
	\label{reg:scelta_pivot}
	Un buon accordo pivot deve soddisfare questi requisiti:
	
	\begin{enumerate}
		\item \textbf{Esistere in entrambe le tonalità} con le stesse note (o con alterazioni compatibili)
		\item \textbf{Avere una funzione forte nella nuova tonalità:} preferire II, IV o V grado della nuova tonalità piuttosto che III o VI
		\item \textbf{Permettere una condotta delle voci fluida} dalla tonalità vecchia a quella nuova
		\item \textbf{Non creare ambiguità premature:} il pivot deve essere chiaramente percepito come parte della tonalità vecchia finché non viene reinterpretato
	\end{enumerate}
\end{regola}

\subsection{Confermare la nuova tonalità}
\label{subsec:conferma_tonalita}

\begin{regola}[Importanza della cadenza di conferma]
	\label{reg:conferma_cadenza}
	Una modulazione non è completa senza una cadenza convincente nella nuova tonalità. Le opzioni più efficaci sono:
	
	\begin{itemize}
		\item \textbf{Cadenza perfetta V--I:} La più forte e inequivocabile
		\item \textbf{Cadenza composta II--V--I o IV--V--I:} Rafforza ulteriormente la nuova tonalità
		\item \textbf{Ripetizione della cadenza:} In caso di dubbio, ripetere la progressione cadenzale elimina ogni ambiguità
	\end{itemize}
	
	Evitare di interrompere la nuova tonalità subito dopo il pivot, senza dare tempo alla nuova tonica di stabilizzarsi.
\end{regola}

\subsection{Gestire le alterazioni}
\label{subsec:gestione_alterazioni}

\begin{regola}[Introduzione delle alterazioni]
	\label{reg:introduzione_alterazioni}
	Quando si modula a una tonalità con un diesis o un bemolle in più:
	
	\begin{itemize}
		\item \textbf{Introdurre l'alterazione gradualmente:} La nuova alterazione appare tipicamente nella sensibile dell'accordo di dominante della nuova tonalità
		\item \textbf{Evitare salti cromatici bruschi:} Se possibile, far apparire l'alterazione in una voce interna piuttosto che nel soprano
		\item \textbf{Preparare l'orecchio:} L'accordo pivot dovrebbe già contenere alcune delle note ``nuove'' della scala di arrivo
	\end{itemize}
\end{regola}

\begin{esempio}[Gestione del \Fadiesis  nella modulazione Do M $\rightarrow$ Sol M]
	\label{ex:gestione_fad}

	\musica{
		\partilate
		{ \key do \major \time 4/4 }
		{ do'2 do'2 | fad'2 sol'2 | sol'1 }
		{ sol2 sol2 | do'2 si2 | si1 }
		{ mi2 mi2 | la2 sol2 | sol1 }
		{ do2 do2 | re2 sol,2 | sol,1 }
		\lyricmode { 
			\markup{"Do M:"} "I (pivot)"2 "I"2 | 
			\markup{"Sol M:"} "IV"2 "V"2 | 
			"I"1 
		}
		\figuremode { s2 s2 | s2 s2 | s1 }
	}
	Il \Fadiesis  (sensibile di Sol Maggiore) appare solo nell'accordo di dominante (Re--\Fadiesis --La--Do), rendendo la modulazione fluida e naturale.
\end{esempio}

\section{Esempi completi di modulazione}
\label{sec:esempi_completi}

Presentiamo ora alcuni esempi completi che illustrano le tre fasi del processo modulante in diversi contesti.

\begin{esempio}[Modulazione completa: La minore $\rightarrow$ Do Maggiore]
	\label{ex:mod_completa_lam_dom}
	Affermazione in La m, pivot attraverso Fa M (VI di La m = IV di Do M), conferma in Do M.

	\musica{
		\partilate
		{ \key la \minor \time 4/4 }
		{ la'2 do''2 | do''2 re''2 | mi''2 do''2 | do''1 }
		{ mi'2 mi'2 | fa'2 fa'2 | do'2 si2 | do'1 }
		{ do'2 do'2 | la2 la2 | sol2 sol2 | mi1 }
		{ la,2 la,2 | fa,2 fa,2 | do2 sol,2 | do1 }
		\lyricmode { 
			\markup{"La m:"} "I"2 "III"2 | 
			"VI (pivot)"2 "VI"2 | 
			\markup{"Do M:"} "IV"2 "V"2 | 
			"I"1 
		}
		\figuremode { s2 s2 | s2 s2 | s2 s2 | s1 }
	}
	La tonalità di partenza (La minore) viene affermata con la progressione I--III. L'accordo di Fa Maggiore funge da pivot, inizialmente percepito come VI di La m, poi reinterpretato come IV di Do M. La cadenza IV--V--I conferma solidamente Do Maggiore.
\end{esempio}

\begin{esempio}[Modulazione con settima: Do M $\rightarrow$ Mi minore]
	\label{ex:mod_settima_do_mim}
	Uso del II$\oslash^7$ di Mi minore come pivot.

	\musica{
		\partilate
		{ \key do \major \time 4/4 }
		{ do'2 do'2 | do'2 si2 | do'2 si2 | si1 }
		{ sol2 la2 | la2 fad2 | sold2 fad2 | mi1 }
		{ mi2 fad2 | re2 re2 | mi2 re2 | si,1 }
		{ do2 re,2 | fad,2 si,,2 | mi,2 si,,2 | mi,,1 }
		\lyricmode { 
			\markup{"Do M:"} "I"2 "II"2 | 
			\markup{"Mi m:"} "II$\oslash^7$"2 "V"2 | 
			"I"2 "V"2 | 
			"I"1 
		}
		\figuremode { s2 s2 | <6 5>2 s2 | s2 s2 | s1 }
	}
\end{esempio}

\section{Regole pratiche per la modulazione rapida}
\label{sec:regole_pratiche_modulazione}

Esistono alcuni procedimenti cromatici ricorrenti che permettono di modulare rapidamente verso le tonalità vicine partendo da una tonalità maggiore.

\begin{regola}[Modulazione al IV grado tramite la sensibile abbassata]
	\label{reg:mod_iv_sensibile_abbassata}
	In una tonalità maggiore, abbassando di un semitono la \textbf{sensibile} (VII grado), si ottiene una nota che funge da \textbf{settima} dell'accordo costruito sulla tonica. Tale accordo (\textbf{I7}) assume funzione di dominante della sottodominante, permettendo una modulazione immediata al \textbf{IV grado}.
	
	\textit{Esempio in Do Maggiore: il Si viene abbassato a \Sibemolle , trasformando l'accordo di Do maggiore in Do7, che risolve su Fa maggiore.}
\end{regola}

\begin{regola}[Modulazione al V grado tramite il IV grado innalzato]
	\label{reg:mod_v_quarto_innalzato}
	Per modulare al \textbf{V grado} (dominante), si agisce sul \textbf{IV grado} della tonalità di partenza innalzandolo di un semitono. Questa nota diventa la \textbf{nuova sensibile} del V grado, trasformando l'accordo sul II grado (o sul IV) in un accordo di dominante della dominante (\textbf{V/V}).
	
	\textit{Esempio in Do Maggiore: il Fa viene innalzato a \Fadiesis , diventando la sensibile di Sol maggiore e preparando l'arrivo della nuova tonica.}
\end{regola}

\begin{regola}[Modulazione al VI grado tramite il V grado innalzato]
	\label{reg:mod_vi_quinto_innalzato}
	Per modulare al \textbf{VI grado} (relativo minore della tonica), si innalza di un semitono il \textbf{V grado} della tonalità di partenza. Questa nota funge da \textbf{nuova sensibile} del VI grado, creando un accordo di dominante (\textbf{V/VI}) che risolve sulla nuova tonica minore.
	
	\textit{Esempio in Do Maggiore: il Sol viene innalzato a \Soldiesis , diventando la sensibile di La minore.}
\end{regola}

\begin{regola}[Modulazione ai relativi minori degli altri gradi]
	\label{reg:mod_altri_minori}
	Analogamente, è possibile modulare ai relativi minori degli altri gradi principali alterando i gradi della scala per ottenere la nuova sensibile:
	
	\begin{itemize}
		\item \textbf{Al III grado (relativo del V):} Si innalzano il \textbf{II} e il \textbf{IV grado} (es. in Do M: \Fadiesis  e \Rediesis  per modulare a Mi minore).
		\item \textbf{Al II grado (relativo del IV):} Si innalza il \textbf{I grado} (es. in Do M: \Dodiesis  per modulare a Re minore).
	\end{itemize}
\end{regola}